\documentclass[12pt, a4paper]{article}

% --------------------------
% Encoding & Language
% --------------------------
\usepackage[utf8]{inputenc}
\usepackage[T1]{fontenc}
\usepackage[english]{babel}

% --------------------------
% Page Layout
% --------------------------
\usepackage[
  top=2.5cm,
  bottom=2.5cm,
  left=3cm,
  right=2.5cm
]{geometry}

% --------------------------
% Typography
% --------------------------
\usepackage{microtype}        % Better text justification
\usepackage{setspace}         % Line spacing control
\onehalfspacing               % 1.5x line spacing

% --------------------------
% Math
% --------------------------
\usepackage{amsmath}
\usepackage{amssymb}

% --------------------------
% Graphics & Colors
% --------------------------
\usepackage{graphicx}
\usepackage[dvipsnames]{xcolor}

% --------------------------
% Hyperlinks (keep as last package)
% --------------------------
\usepackage[
  colorlinks=true,
  linkcolor=blue,
  urlcolor=blue,
  citecolor=blue
]{hyperref}

% --------------------------
% Document Metadata
% --------------------------
\title{Generic \LaTeX\ Document}
\author{Your Name}
\date{\today}

% ==========================
% BEGIN DOCUMENT
% ==========================
\begin{document}

\maketitle
\tableofcontents
\newpage

% --------------------------
\section{Introduction}
% --------------------------

This is a generic \LaTeX\ template. Feel free to adapt it to your needs — 
whether you're writing a report, an article, or just experimenting with 
formatting.

You can write normal paragraphs by just typing text. A blank line between
paragraphs starts a new one.

% --------------------------
\section{Text Formatting}
% --------------------------

Here are some common text styles:

\begin{itemize}
  \item \textbf{Bold text} with \texttt{\textbackslash textbf\{\}}
  \item \textit{Italic text} with \texttt{\textbackslash textit\{\}}
  \item \underline{Underlined text} with \texttt{\textbackslash underline\{\}}
  \item \texttt{Monospaced text} with \texttt{\textbackslash texttt\{\}}
\end{itemize}

% --------------------------
\section{Math}
% --------------------------

Inline math uses dollar signs: $E = mc^2$.

Display math uses the \texttt{equation} environment:

\begin{equation}
  \int_{-\infty}^{\infty} e^{-x^2} \, dx = \sqrt{\pi}
\end{equation}

Or use \texttt{align} for multiple equations:

\begin{align}
  f(x) &= x^2 + 3x + 2 \\
       &= (x+1)(x+2)
\end{align}

% --------------------------
\section{Lists}
% --------------------------

Unordered list:
\begin{itemize}
  \item First item
  \item Second item
  \item Third item
\end{itemize}

Ordered list:
\begin{enumerate}
  \item Step one
  \item Step two
  \item Step three
\end{enumerate}

% --------------------------
\section{Tables}
% --------------------------

\begin{table}[h!]
  \centering
  \caption{A simple example table}
  \begin{tabular}{|c|c|c|}
    \hline
    \textbf{Column A} & \textbf{Column B} & \textbf{Column C} \\
    \hline
    Value 1 & Value 2 & Value 3 \\
    Value 4 & Value 5 & Value 6 \\
    \hline
  \end{tabular}
  \label{tab:example}
\end{table}

% --------------------------
\section{Including Figures}
% --------------------------

To include an image, place the file (e.g., \texttt{my-image.png}) in the 
same folder and use:

\begin{verbatim}
\begin{figure}[h!]
  \centering
  \includegraphics[width=0.7\textwidth]{my-image}
  \caption{Your caption here}
  \label{fig:example}
\end{figure}
\end{verbatim}

% --------------------------
\section{Cross-References}
% --------------------------

You can reference any labeled element, like Table~\ref{tab:example},
using \texttt{\textbackslash ref\{\}} and \texttt{\textbackslash label\{\}}.
Hyperlinks work too: \href{https://www.latex-project.org}{LaTeX Project website}.

% ==========================
\end{document}